\documentclass[11pt,letterpaper]{article}
\usepackage[T1]{fontenc}
\usepackage[utf8]{inputenc}
\usepackage{mathtools,amsthm}
\usepackage{newpxtext,newpxmath}
\usepackage[margin=1in,headheight=15pt]{geometry}
\usepackage{microtype}
\usepackage{booktabs,tabularx,array,longtable}
\usepackage{xcolor}
\usepackage[most]{tcolorbox}
\usepackage{enumitem}
\usepackage{fancyhdr}
\usepackage{titlesec}
\usepackage{xurl}
\usepackage{hyperref}
\usepackage[nameinlink,noabbrev]{cleveref}
\definecolor{oliveink}{HTML}{4F6041}
\definecolor{softolive}{HTML}{F1F3EB}
\definecolor{warmgray}{HTML}{66645E}
\definecolor{rulegray}{HTML}{D5D8CA}
\hypersetup{colorlinks=true,linkcolor=oliveink,citecolor=oliveink,urlcolor=oliveink,
  pdftitle={A countable genetic cut across the finite--infinite gap},
  pdfauthor={Research note prepared for Vladimir Reshetnikov},
  pdfsubject={Explicit counterexamples concerning genetic surreal functions}}
\titleformat{\section}{\Large\bfseries\color{oliveink}}{\thesection}{0.65em}{}
\titleformat{\subsection}{\large\bfseries}{\thesubsection}{0.65em}{}
\pagestyle{fancy}
\fancyhf{}
\fancyhead[L]{\small\color{warmgray}Genetic cuts and surreal primitives}
\fancyhead[R]{\small\color{warmgray}Research note \;|\; 20 September 2026}
\fancyfoot[C]{\thepage}
\renewcommand{\headrulewidth}{0.3pt}
\setlength{\parindent}{1.2em}
\setlength{\parskip}{0.22em}
\setlength{\emergencystretch}{2em}
\setcounter{tocdepth}{1}
\numberwithin{equation}{section}
\newtheorem{theorem}{Theorem}[section]
\newtheorem{lemma}[theorem]{Lemma}
\newtheorem{proposition}[theorem]{Proposition}
\newtheorem{corollary}[theorem]{Corollary}
\theoremstyle{definition}
\newtheorem{definition}[theorem]{Definition}
\newtheorem{example}[theorem]{Example}
\newtheorem{question}[theorem]{Question}
\theoremstyle{remark}
\newtheorem{remark}[theorem]{Remark}
\newcommand{\No}{\mathbf{No}}
\newcommand{\On}{\mathbf{On}}
\newcommand{\Off}{\mathbf{Off}}
\newcommand{\N}{\mathbb{N}}
\newcommand{\Z}{\mathbb{Z}}
\newcommand{\Q}{\mathbb{Q}}
\newcommand{\R}{\mathbb{R}}
\newcommand{\Fin}{\mathcal{O}}
\newcommand{\U}{\mathcal{U}}
\newcommand{\D}{\mathcal{D}}
\newcommand{\cut}[2]{\left\{\,#1\;\middle|\;#2\,\right\}}
\newcommand{\ind}{\mathbf{1}}
\newcommand{\dd}{\mathrm{d}}
\newcommand{\sgn}{\operatorname{sgn}}
\newcommand{\id}{\operatorname{id}}
\newcommand{\bday}{\operatorname{b}}
\newtcolorbox{resultbox}{colback=softolive,colframe=rulegray,boxrule=0.6pt,
  arc=1.5mm,left=3mm,right=3mm,top=2mm,bottom=2mm,breakable}
\newcommand{\source}[2]{\cite[#1]{#2}}

\begin{document}
\thispagestyle{empty}
\begin{center}
{\small\sffamily\color{oliveink} SURREAL ANALYSIS \quad / \quad EXPLICIT CONSTRUCTION}\par
\vspace{0.75cm}
{\LARGE\bfseries A countable genetic cut across\\[0.12em] the finite--infinite gap}\par
\vspace{0.4cm}
{\large Counterexamples concerning uniqueness of primitives\\
and the sup property of entire genetic functions}\par
\vspace{0.5cm}
{\normalsize Research note prepared for Vladimir Reshetnikov}\par
\vspace{0.15cm}
{\small 20 September 2026}
\end{center}
\vspace{0.35cm}
\begin{abstract}
Put $q(t)=t/(1+t^2)$ and form the countable Conway cut
\[
 T(x)=\cut{q(n-x):n\in\N}{},\qquad B(x)=1-T(x).
\]
We prove that $B$ is the indicator of the positive infinite surreal numbers.
The option functions are globally defined rational functions, so this is an
admissible genetic construction under the set-sized option rules in
Rubinstein-Salzedo--Swaminathan and in Berenbeim's later formulation; the
latter's global cofinality requirement is checked directly. Yet $B$ is locally
constant everywhere and nonconstant globally. Consequently $x$ and $x-B(x)$
are genetic primitives of $1$ that agree on every real argument but do not
differ by a constant. The second is an increasing bijection of $\No$.
A bounded sublevel class of $B$ has no surreal supremum, giving a second
counterexample. We extend the construction to arbitrary nonempty sets of
surreals without a largest element, construct ordinal-indexed independent
families in the kernel of pointwise differentiation, and explain precisely
why real-variable uniqueness arguments do not apply. The accompanying
standard-library Python program performs exact finite checks, not a formal
verification of the infinitary proofs.
\end{abstract}
\begin{resultbox}
\textbf{Status and provenance.}
The targets are Conjecture~53 of \cite{RSS} and Chapter~8,
Conjecture~3 of \cite{Berenbeim}. The piecewise primitive itself already
appears in \cite{RSS}; the contribution here is the explicit genetic formula
and the deductions from it. The arguments below give a negative answer
under the published option conventions. No later resolution was located in
the searches recorded with this note, but those searches do not establish
priority or exclude unpublished work. The proof has not been independently
refereed or checked in a proof assistant.
\end{resultbox}
\newpage
\tableofcontents
\vspace{0.6cm}
\noindent\textbf{Reading guide.}
Sections~\ref{sec:prelim}--\ref{sec:primitive} contain the complete central
argument. Section~\ref{sec:admissibility} is essential: a nonconstant
zero-derivative function alone would not answer the genetic-function
conjecture. Section~\ref{sec:sup} treats the second target.
The remaining sections give extensions, failure analyses, and a precise
account of what the computational checks establish.

\section{The published questions and the intended scope}\label{sec:targets}
A primitive of a function $f$ is a function $F$ with $F'=f$. The question
studied here asks whether restricting both $f$ and its primitives to
\emph{genetic} surreal functions restores uniqueness up to a constant.
This is the substance of Rubinstein-Salzedo and Swaminathan's Conjecture~53
\source{p.~37}{RSS}. Berenbeim's later Conjecture~3 instead asks whether
every entire genetic function has the \emph{sup property}
\source{pp.~225--226}{Berenbeim}.

The term \emph{genetic} refers to specified constructions of functions
using Conway cuts, previously available functions, and recursion on simpler
arguments. It does not mean merely ``definable,'' and it does not by itself
mean analytic, continuous, or finitely presented. The distinction matters:
a direct piecewise definition is not, without further work, a certificate
of membership in a genetic class.

Our main construction requires only a particularly small part of the
published machinery: rational functions with no poles, a countable fixed
family of options, a single nonrecursive cut, and ordinary arithmetic on
the resulting functions. We state the operational requirements explicitly
in Section~\ref{sec:admissibility} instead of assuming that every reasonable
piecewise function is genetic.

There are two separate issues of validation. The mathematical issue is
whether the displayed construction has the stated properties under the
specified rules. That is addressed by proofs. The historical issue is
whether this particular construction or an equivalent argument is already
known. The searches supporting this note locate the published conjectures
and their relevant definitions; they do not settle historical priority.

\section{Cuts, infinitude, and the derivative in use}\label{sec:prelim}
\subsection{The small amount of surreal theory needed}
The surreal numbers form a proper-class ordered field. If $L$ and $R$ are
\emph{sets} of surreal numbers and every element of $L$ is less than every
element of $R$, then $\cut{L}{R}$ is the simplest surreal strictly between
them. Empty option sets are allowed. We use the standard simplicity rule,
in particular $0=\cut{}{}$ and $1=\cut{0}{}$; these conventions are recalled
in \source{Section~2.1}{RSS}.

The braces in $\cut{L}{R}$ do \emph{not} denote an ordinary real supremum.
For example, a set of negative left options may approach $0$ in various
senses, but if there are no right options, $0$ itself is already the simplest
number above all of them. Conversely, one nonnegative left option excludes
$0$ even if that option is an infinitesimal.

All algebra below is \emph{surreal field algebra}. In particular, $\omega-1$
is meaningful in $\No$ and is not an ordinal predecessor of $\omega$.
Whenever ordinal exponents are used later, $\omega^\alpha$ denotes the
usual embedded ordinal monomial. Neither noncommutative ordinal addition
nor ordinal subtraction is being used in the rational formulas.

\begin{definition}\label{def:gap}
Throughout, $\N=\{0,1,2,\ldots\}$ denotes the ordinary, finite natural
numbers. Define
\begin{align}
 \U&=\{x\in\No:(\forall n\in\N)\ x>n\},\\
 \D&=\{x\in\No:(\exists n\in\N)\ x\le n\},\\
 \Fin&=\{x\in\No:(\exists n\in\N)\ |x|\le n\}.
\end{align}
Thus $\U$ is the class of positive infinite numbers, $\Fin$ is the class of
finite numbers, and $\D=\No\setminus\U$ contains both the finite numbers
and all negative infinite numbers.
\end{definition}

There is no surreal number sitting at the boundary between $\D$ and $\U$.
The left class has no greatest element: from $x\le n$ we get $n+1>x$ with
$n+1\in\D$. The right class has no least element: if $u\in\U$, then
$0<u/2<u$ and $u/2\in\U$, since $u>2n$ for every natural $n$.
We shall refer to this boundary as the \emph{finite--infinite gap}. The
word ``finite'' in this name is not a claim that every member of $\D$ is
finite.

A useful stronger separation is
\begin{equation}\label{eq:crossgap}
 d\in\D,\ u\in\U\quad\Longrightarrow\quad u-d\in\U.
\end{equation}
Indeed, choose $n$ with $d\le n$. For any $m\in\N$ we have
$u>m+n$, so $u-d\ge u-n>m$.

\subsection{Local differentiation, not a derivation of field elements}
For a function on a neighborhood of $x$, use
\begin{equation}\label{eq:derivative}
 F'(x)=\lim_{h\to0}\frac{F(x+h)-F(x)}{h}.
\end{equation}
An order-field $\varepsilon$--$\delta$ limit quantifies over all positive
surreal $\varepsilon$ and $\delta$, not only positive real ones. This is
the pointwise derivative at issue in \source{Section~5}{RSS}. Our difference
quotients will be exactly constant on a punctured neighborhood, so there is
no subtle convergence claim to establish.

This operator acts on \emph{functions}. It should not be confused with a
chosen derivation $\partial:\No\to\No$ on the surreal field itself.
The assertion $F'(x)=1$ concerns the variation of $F$ with its argument;
it is not an assertion about the value of a field derivation at the number
$F(x)$. No result about a distinguished field derivation is being challenged.

\section{A countable cut that detects positive infinity}\label{sec:cut}
\begin{lemma}[A bounded sign-preserving rational function]\label{lem:q}
In every ordered field, the function
\begin{equation}\label{eq:q}
 q(t)=\frac{t}{1+t^2}
\end{equation}
is everywhere defined, has the same sign as $t$, and satisfies
$-1<q(t)<1$. In fact, $|q(t)|\le\tfrac12$.
\end{lemma}
\begin{proof}
The denominator is positive. For the strict bounds, the identities
\[
 1+t^2-t=(t-\tfrac12)^2+\tfrac34>0,\qquad
 1+t^2+t=(t+\tfrac12)^2+\tfrac34>0
\]
give $q(t)<1$ and $q(t)>-1$. The sharper bound follows from
$(|t|-1)^2\ge0$, equivalently $2|t|\le1+t^2$.
\end{proof}

\begin{lemma}[A binary cut test]\label{lem:binary}
Let $L$ be a set of surreal numbers, all strictly less than $1$. If every
member of $L$ is negative, then $\cut{L}{}=0$. If $L$ contains a nonnegative
member, then $\cut{L}{}=1$.
\end{lemma}
\begin{proof}
In the first case $0$ lies above every left option and is the simplest
surreal of all. In the second case $0$ is excluded, whereas $1$ lies above
all left options. The only other number born on the day when $1$ is born
is $-1$, and it is excluded by the nonnegative option. Thus $1$ is the
simplest admissible number. The first case also covers $L=\varnothing$.
\end{proof}

\begin{definition}\label{def:T}
For every $x\in\No$, set
\begin{equation}\label{eq:T}
 T(x)=\cut{\frac{n-x}{1+(n-x)^2}:n\in\N}{},\qquad B(x)=1-T(x).
\end{equation}
\end{definition}

For each fixed $x$, the left side of this cut is a countable set of surreal
numbers. There are no right options, so the order condition is vacuous.
The definition therefore produces a surreal number at every input. The
natural-number quantifier is part of the option indexing; it is not a
limit operation and it is not replaced by an ordinal-length summation.

\begin{theorem}[Exact evaluation]\label{thm:evaluation}
The functions in \cref{def:T} satisfy
\begin{equation}\label{eq:values}
 T(x)=
 \begin{cases}1,&x\in\D,\\0,&x\in\U,\end{cases}
 \qquad
 B(x)=
 \begin{cases}0,&x\in\D,\\1,&x\in\U.\end{cases}
\end{equation}
\end{theorem}
\begin{proof}
Every left option lies in $[-\tfrac12,\tfrac12]$, by \cref{lem:q}.
If $x\in\D$, choose $n\in\N$ with $x\le n$. Then $n-x\ge0$, so the
corresponding left option is nonnegative. The second case of
\cref{lem:binary} gives $T(x)=1$.

If $x\in\U$, then $n-x<0$ for every $n\in\N$. All left options are
negative, so the first case gives $T(x)=0$. Subtracting from $1$ gives the
formula for $B$.
\end{proof}

\begin{example}
The distinction remains exact for infinitesimals and infinite numbers:
\[
\begin{array}{c|cccccccc}
x&-\omega&0&1/\omega&1&1+1/\omega&\omega/2&\omega-1&\omega^2\\ \hline
T(x)&1&1&1&1&1&0&0&0\\
B(x)&0&0&0&0&0&1&1&1
\end{array}
\]
For example, the $n=2$ option is positive at $x=1+1/\omega$, while every
finite-$n$ option is negative at $x=\omega/2$. No approximation of a sign is
being made.
\end{example}

\section{Why the construction is genetic}\label{sec:admissibility}
\subsection{The admissibility certificate}
We isolate the exact constructor being used. Begin with the rational
function $q$, take a fixed set of affine translates of it, and use their
values as left options in a Conway cut. There is no reference to the
function currently being defined, nor to a left or right option of the
input. Such a formula is a \emph{nonrecursive instance} of a genetic
constructor, not an additional operation being postulated here.

In the original construction the admissible option expressions include
$c_1h(c_2x+c_3)+c_4$, where $h$ is already genetic and the $c_i$ may be
surreal constants. The chosen option families are set-sized. Arithmetic
closure is also available \source{pp.~9--10}{RSS}. For our $n$th option use
\[
 h=q,\qquad c_1=1,\quad c_2=-1,\quad c_3=n,\quad c_4=0.
\]
Thus $q(n-x)$ is an allowed expression before $T$ is introduced. Its
membership does not depend on knowing whether $x$ is infinite.

The original paper explicitly accepts $q$ and the one-option cut
$\cut{q(x)}{}$ as genetic \source{p.~11}{RSS}. This matters because it
establishes that an option expression depending on $x$ but on no options
of $x$ is admissible. We are using the same type of expression with a
countable option set instead of a singleton.

\begin{proposition}\label{prop:genetic}
$T$, $B$, $\id$, and $\id-B$ are genetic under the constructor just
specified. The option formula for $T$ is independent of the representation
chosen for $x$.
\end{proposition}
\begin{proof}
The set of option expressions
\[
 \mathscr L=\{x\mapsto q(n-x):n\in\N\},\qquad
 \mathscr R=\varnothing
\]
is countable and consists of allowed previously constructed expressions.
For any $x$, all options have defined surreal values. No left--right
conflict is possible. The genetic adjunction therefore defines $T$.
Polynomial and arithmetic closure gives $B=1-T$, $\id$, and $\id-B$.

Changing a presentation $x=\cut{L_x}{R_x}$ changes none of the numerical
expressions $q(n-x)$. Hence it changes neither option set nor cut value.
This is uniformity with respect to the presentation of $x$.
\end{proof}

\subsection{The stronger cofinality check}
Berenbeim's Definition~55 uses a positive-denominator localization and an
additional global cofinality condition \source{pp.~126--128}{Berenbeim}.
For a left option expression $\ell(v,w;x)$, the relevant requirement is
that $\ell(y,z;x)<T(x)$ whenever $y<x<z$; a dual inequality applies to
right options. Here
\[
 \ell_n(v,w;x)=q(n-x)
\]
ignores both $v$ and $w$. If $x\in\D$, then
$\ell_n(y,z;x)\le\tfrac12<1=T(x)$. If $x\in\U$, then
$\ell_n(y,z;x)<0=T(x)$. The right-option condition is vacuous.

The denominator $1+(n-x)^2$ is strictly positive on the entire domain, so
it is also compatible with the localization. Thus the certificate checks
the extra condition directly, rather than deducing it from an informal
assertion that the function ``looks uniform.'' The thesis also explicitly
treats the one-option rational step construction as entire and uniform
\source{pp.~133--134}{Berenbeim}.

\begin{remark}[What ``entire'' means here]
In this setting ``entire'' concerns a genetic function defined at every
surreal argument, without a pole or an excluded input. It is not the
complex-analytic notion of a globally holomorphic function. Formula
\eqref{eq:T} is entire in the former sense.
\end{remark}

\subsection{A notational issue that must not be hidden}
Genetic formulas often display unions over $x^L$ and $x^R$ even when a
particular option term uses neither variable. Such a term remains present
when one or both input option sets are empty. This convention is necessary
already for the accepted one-option step function at $x=0$. Reading a
schematic union as a literal Cartesian-product index that deletes every
independent term when $R_x=\varnothing$ would not reproduce that example.
Fornasiero makes the convention about unused option variables explicit
\source{Example~2.5}{Fornasiero}.

Our proof therefore uses the usual option-expression semantics exhibited
by the sources themselves. It does not exploit undefined terms, hidden
conditional inclusion, class-sized options, or a choice of an unusual
presentation of $x$. A materially narrower notion, for example one
allowing only finitely many option expressions, would be a different
question; this note does not settle it.

\section{Nonunique primitives, even with the same real values}\label{sec:primitive}
\begin{lemma}[Invariance under finite translation]\label{lem:translation}
If $c\in\Fin$, then
\[
 x\in\U\quad\Longleftrightarrow\quad x+c\in\U.
\]
Consequently $B(x+c)=B(x)$ and $T(x+c)=T(x)$.
\end{lemma}
\begin{proof}
Choose $k\in\N$ with $|c|\le k$. If $x\in\U$, then for every $n\in\N$,
$x>n+k$, and therefore $x+c\ge x-k>n$. Thus $x+c\in\U$.
The reverse implication follows by applying the same argument to $-c$.
\end{proof}

\begin{theorem}[The kernel is not just the constants]\label{thm:zero}
$B$ and $T$ are locally constant at every surreal argument. More precisely,
\begin{equation}\label{eq:radiusone}
 |h|<1\quad\Longrightarrow\quad B(x+h)=B(x)
 \quad\text{for every }x\in\No.
\end{equation}
Hence $B'(x)=T'(x)=0$ for all $x$, although neither function is constant.
All their higher pointwise derivatives exist and vanish.
\end{theorem}
\begin{proof}
Every $h$ with $|h|<1$ is finite, so \cref{lem:translation} proves
\eqref{eq:radiusone}. For $0<|h|<1$ the difference quotient for $B$ is
exactly $0$. Given any surreal $\varepsilon>0$, choosing $\delta=1$ makes
its distance from $0$ less than $\varepsilon$. The same applies to $T$.

But $B(0)=0$ and $B(\omega)=1$, so $B$ is not constant, and neither is
$T=1-B$. The first derivative is the zero function; repeated differentiation
therefore gives zero at every higher order.
\end{proof}

The radius $1$ in \eqref{eq:radiusone} is independent of $x$ and of the
requested accuracy. The apparent paradox is that no finite chain of
steps of size less than $1$ can carry a finite number to a positive infinite
number. A real proof that connects two points by finitely many such steps
uses the Archimedean property, which is unavailable here.

\begin{theorem}[Counterexample to genetic primitive uniqueness]\label{thm:primitive}
Define, on all of $\No$,
\begin{equation}\label{eq:F}
 F_0(x)=x,\qquad F_1(x)=x-B(x)=x-1+T(x).
\end{equation}
Then $F_0$ and $F_1$ are genetic, $F_0'=F_1'=1$ everywhere, and
$F_0(r)=F_1(r)=r$ for every real $r$. Nevertheless $F_1-F_0$ is not constant.
The same failure occurs on the proper bounded open interval $(0,\omega)$.
\end{theorem}
\begin{proof}
Geneticity is \cref{prop:genetic}. For $0<|h|<1$,
\[
 \frac{F_1(x+h)-F_1(x)}h
 =\frac{h-B(x+h)+B(x)}h=1.
\]
Thus its derivative is $1$, and the identity function has the same
derivative. Every real number belongs to $\D$, so $B(r)=0$ and the real
values agree. The difference is $-B$, which takes both $0$ and $-1$.

On $(0,\omega)$ use the two points $1$ and $\omega/2$. Both belong to the
interval, and the differences there are $0$ and $-1$, respectively.
For differentiation within this interval, shrink the radius at $x$ to
$\min\{1,x/2,(\omega-x)/2\}>0$. This leaves the same quotient calculation
valid on a neighborhood contained in the domain.
\end{proof}

The piecewise description of $F_1$ is precisely
\[
 F_1(x)=\begin{cases}x,&x\in\D,\\x-1,&x\in\U.\end{cases}
\]
This is the function already displayed in equation~(7) immediately before
the conjecture in \source{p.~37}{RSS}. The present answer does not claim to
have discovered this piecewise primitive. It supplies the genetic
representation that makes it a counterexample to the restriction proposed
there.

\begin{corollary}[Failure of normalized global uniqueness]
The global initial-value problem $F'=1$, $F(0)=0$ has at least two genetic
solutions. Even prescribing $F(r)=r$ simultaneously for every $r\in\R$
does not select one of them. Likewise $G'=0$, $G(0)=0$ has the distinct
genetic solutions $G=0$ and $G=B$.
\end{corollary}
\begin{proof}
Apply \cref{thm:primitive,thm:zero}. No assertion of failure of local
uniqueness on a sufficiently small neighborhood of $0$ is intended:
the displayed solutions agree on every finite argument.
\end{proof}

\section{Failure of the sup property}\label{sec:sup}
For a single-variable function $f$, the sup property requires that for
$a<b$ and $c$ in the field, the infima and suprema of both classes
\[
 \{x:a<x<b,\ f(x)\le c\},\qquad
 \{x:a<x<b,\ f(x)\ge c\}
\]
are field elements or the two end gaps. This is Fornasiero's
Definition~3.2 \cite{Fornasiero}, used in the later conjecture.
It is enough to exhibit one nonempty bounded sublevel class with no
surreal supremum.

\begin{theorem}[Counterexample to the entire-genetic sup conjecture]\label{thm:sup}
The entire genetic function $B$ fails the sup property. In particular,
\begin{equation}\label{eq:E}
 E=\{x\in\No:0<x<\omega,\ B(x)\le\tfrac12\}
\end{equation}
has no supremum in $\No\cup\{\Off,\On\}$.
\end{theorem}
\begin{proof}
By \cref{thm:evaluation}, $E$ is exactly the class of positive finite
surreals. The condition $x<\omega$ is automatic for every such number.
The class contains every positive integer and is bounded above by $\omega$.

If $u$ is any surreal upper bound of $E$, then $u\ge n+1>n$ for every
$n\in\N$, since $n+1\in E$. Hence $u\in\U$. Conversely, every $u\in\U$
is an upper bound of $E$: for $e\in E$, choose $n\in\N$ with $e\le n<u$.
Thus the class of surreal upper bounds of $E$ is exactly $\U$.

This class has no least member, because $u/2\in\U$ and $u/2<u$ for every
$u\in\U$. There is therefore no surreal least upper bound. The bottom end
gap cannot be an upper bound of a nonempty positive class, and the top end
gap cannot be least because $\omega$ is already an upper bound. This proves
the claim.
\end{proof}

For completeness, the complementary superlevel class in the same interval,
\[
 E^+=\{x:0<x<\omega,\ B(x)\ge\tfrac12\}=\U\cap(0,\omega),
\]
has no surreal infimum. Every $d\in\D$ is a lower bound. If a lower bound
$u$ belonged to $\U$, then either $u<\omega$, in which case
$u/2\in E^+$ lies below it, or $u\ge\omega$, in which case
$\omega/2\in E^+$ lies below it. Thus all lower bounds lie in $\D$, and
$\D$ has no greatest member.

Notice the precise obstruction: boundedness supplies many upper bounds,
not a \emph{least} upper bound. Being defined everywhere and having no
rational poles does not remove an unrealized Dedekind boundary of a
sublevel class.

\section{Every set-generated endpoint-free cut gives a variant}\label{sec:general}
The use of $\N$ is not essential for the evaluation mechanism. What
matters for differentiability is whether the lower indexing set has a
last element.

\begin{definition}
For a nonempty set $A\subset\No$, define
\begin{align}
 T_A(x)&=\cut{q(a-x):a\in A}{},& B_A(x)&=1-T_A(x),\\
 \U_A&=\{x:(\forall a\in A)\ a<x\},&
 \D_A&=\{x:(\exists a\in A)\ x\le a\}.
\end{align}
The set $A$ is fixed before the function is defined.
\end{definition}

\begin{theorem}[General gap-indicator construction]\label{thm:general}
For every nonempty set $A\subset\No$, $T_A$ and $B_A$ are entire uniform
genetic functions under the same option rules, with
$T_A=\ind_{\D_A}$ and $B_A=\ind_{\U_A}$. Both values occur.
If $A$ has no largest element, both functions are locally constant
everywhere and have zero derivative everywhere. If $A$ has a largest
element $m$, they are discontinuous at $m$ and are not differentiable
there.
\end{theorem}
\begin{proof}
The option expressions are a set of affine translates of $q$ with surreal
constants. They ignore all options of $x$, and their denominators are
positive. The construction and cofinality checks of
Section~\ref{sec:admissibility} apply with $a$ in place of $n$.
The sign argument of \cref{thm:evaluation} proves the stated values.
Any $a\in A$ lies in $\D_A$, while $\cut{A}{}$ lies in $\U_A$, so the
function is nonconstant.

Suppose first that $A$ has no largest element. If $x\in\D_A$, choose
$a\in A$ with $x\le a$ and then $a'\in A$ with $a'>a$. Put
$\delta=(a'-x)/2$. For $|h|<\delta$ we have $x+h<a'$, so
$x+h\in\D_A$.

If $x\in\U_A$, form the legitimate set-sized cut
$y=\cut{A}{x}$. It satisfies $a<y<x$ for every $a\in A$. With
$\delta=(x-y)/2$, the inequality $|h|<\delta$ gives $x+h>y>a$ for
all $a\in A$, hence $x+h\in\U_A$. Thus $B_A$ and $T_A$ are locally
constant and their difference quotients vanish near each point.

Suppose instead that $m=\max A$. Then $B_A(x)=0$ for $x\le m$ and
$B_A(x)=1$ for $x>m$. In every neighborhood of $m$ there are points of
both types, so continuity fails with tolerance $\tfrac12$.
For a direct derivative argument, if a derivative $L$ at $m$ existed,
choose $h>0$ smaller than both its proposed tolerance-$1$ radius and
$1/(2(|L|+1))$. The quotient equals $1/h>2(|L|+1)$ and hence cannot
be within $1$ of $L$. Thus no derivative exists at $m$.
\end{proof}

A finite nonempty $A$ has a maximum, so it cannot produce an everywhere
zero-derivative nonconstant function by this particular formula. A
countable set without a maximum already can. This is a minimality
statement about the \emph{displayed schema}, not about all possible
finite genetic definitions.

\begin{proposition}[Cofinality is the invariant]
Suppose $A,C\subset\No$ are nonempty sets and each is weakly cofinal in
the other: every member of either set is at most some member of the other.
Then $T_A=T_C$ and $B_A=B_C$.
\end{proposition}
\begin{proof}
A number is strictly greater than all of $A$ if and only if it is strictly
greater than all of $C$. Indeed, an $a\in A$ can be bounded by $c\in C$,
and the same argument works in reverse. Thus $\U_A=\U_C$, and the
explicit indicator formulas apply.
\end{proof}

The derivatives therefore depend on the endpoint-free cut encoded by
$A$, not on a particular enumeration or on the detailed magnitudes of
the rational options. No real supremum of $A$ needs to exist.

\section{Order automorphisms and an ordinal-indexed kernel}\label{sec:extensions}
\subsection{A group of normalized primitives}
For $c\in\No$, let
\begin{equation}\label{eq:Fc}
 \Phi_c(x)=x+cB(x).
\end{equation}
Every $\Phi_c$ is genetic and has derivative $1$, whatever the size of $c$.
Finite $c$ give a stronger conclusion.

\begin{theorem}\label{thm:auto}
$\Phi_c$ is an increasing bijection $\No\to\No$ exactly when $c$ is finite.
For finite $c,d$,
\begin{equation}\label{eq:group}
 \Phi_c\circ\Phi_d=\Phi_{c+d},\qquad
 \Phi_c^{-1}=\Phi_{-c}.
\end{equation}
\end{theorem}
\begin{proof}
Let $c$ be finite. Translation by $c$ preserves both $\U$ and $\D$ by
\cref{lem:translation}. On $\D$, $\Phi_c$ is the identity; on $\U$, it is
translation by $c$, a bijection of that class. It is strictly increasing
on each class. If $d\in\D$ and $u\in\U$, then $u-d$ is positive infinite
by \eqref{eq:crossgap}, so $(u-d)+c>0$. Thus
$\Phi_c(u)=u+c>d=\Phi_c(d)$, establishing strict increase across the gap.

Since $B(\Phi_d(x))=B(x)$ for finite $d$, substitution gives
\[
 \Phi_c(\Phi_d(x))=x+dB(x)+cB(\Phi_d(x))
 =x+(c+d)B(x).
\]
This proves \eqref{eq:group}, including the inverse formula.

Conversely, if $c$ is positive infinite, then $c$ is not in the range of
$\Phi_c$: points in $\D$ remain in $\D$, while for $x\in\U$,
$x+c>c$. Thus the map is not onto. If $c$ is negative infinite, then
$-c\in\U$ and $\Phi_c(-c)=0=\Phi_c(0)$, so the map is not injective.
These exhaust the nonfinite cases.
\end{proof}

Taking $c=-1$ recovers $F_1$. Thus the central counterexample is not merely
an arbitrary discontinuous rearrangement of values: it is an order
automorphism with exactly the same derivative as the identity.

\subsection{Many independent zero-derivative functions}
For each ordinal $\alpha$, put
\[
 s_\alpha=\omega^\alpha,\qquad
 A_\alpha=\{ns_\alpha:n\in\N\},\qquad B_\alpha=B_{A_\alpha}.
\]
Here $ns_\alpha$ is field multiplication. Because $s_\alpha>0$,
$A_\alpha$ has no maximum, so every $B_\alpha$ is a nonconstant entire
genetic function with zero pointwise derivative.

\begin{theorem}[Ordinal-indexed finite linear independence]\label{thm:independent}
If $\alpha_1<\cdots<\alpha_k$ are ordinals and $c_1,\ldots,c_k\in\No$,
then
\[
 \sum_{i=1}^{k}c_iB_{\alpha_i}=0
 \quad\Longrightarrow\quad c_1=\cdots=c_k=0.
\]
Consequently there are independent families of these kernel elements
indexed by sets of arbitrarily large ordinal order type.
\end{theorem}
\begin{proof}
Evaluate at $x_j=\omega^{\alpha_j+1}$. If $i\le j$, then
$x_j>n\omega^{\alpha_i}$ for every natural $n$, so
$B_{\alpha_i}(x_j)=1$. For $i=j$ this follows from
$\omega^{\alpha_j+1}=\omega^{\alpha_j}\omega$; for earlier $i$ the
scale is no larger. If $i>j$, ordinal order gives
$\alpha_i\ge\alpha_j+1$, hence $x_j\le\omega^{\alpha_i}\in A_{\alpha_i}$,
so $B_{\alpha_i}(x_j)=0$.

Thus the evaluation matrix has entries
\[
 B_{\alpha_i}(x_j)=\begin{cases}1,&i\le j,\\0,&i>j.\end{cases}
\]
The first evaluation gives $c_1=0$, the second gives $c_1+c_2=0$, and
successive subtraction gives every coefficient zero.
\end{proof}

This is a statement about finite linear relations, so no summability
convention for a proper-class family is needed. Each finite linear
combination remains genetic and locally constant: at a fixed point,
intersect finitely many neighborhoods on which its summands are constant.
Adding such a combination to $x$ gives further primitives of $1$.
We do not assign a set-valued dimension to a proper-class function space.

\section{Why familiar uniqueness arguments do not apply}\label{sec:obstruction}
\subsection{The missing point in a mean-value argument}
Over the real numbers, zero derivative on an interval implies constancy
through the mean value theorem. That theorem is not an algebraic
consequence of the ordered-field axioms. Our example makes the failing
implication explicit: $B$ has derivative $0$ everywhere but
$B(\omega)-B(0)=1$.

There is no discontinuity at a hidden surreal ``infinity point.'' No such
point exists. Every actual surreal has a neighborhood entirely on one
side of the gap. In particular, testing differentiability at $\omega$,
$\omega/2$, $\omega-1$, or an arbitrary positive infinite number cannot
locate a jump. The jump belongs to a Dedekind boundary not realized by a
field element.

The failure is visible on the bounded interval $[0,\omega]$. Boundedness
in a non-Archimedean field is not compactness or Dedekind completeness.
A supremum proof of Rolle's theorem cannot simply be imported from real
analysis: the class whose endpoint one needs may have a gap as its
supremum, as in \cref{thm:sup}.

\subsection{Local openness versus the set-union topology}
The original article uses an openness convention built from
\emph{set-sized} unions of intervals whose endpoints are surreal numbers
or the two end gaps \source{Definition~7}{RSS}. This is weaker than
allowing every locally open proper class to count as open. The distinction
can be seen directly here.

Both $\D$ and $\U$ are locally open. Also
\[
 \D=\bigcup_{n\in\N}(-\infty_{\rm end},n)
\]
is a set-sized union of allowed intervals. But $\U$ is not such a union.
To prove this, suppose $\U=\bigcup_{i\in I}(a_i,b_i)$ with $I$ a set.
Every nonempty interval in the union must have its lower endpoint
$a_i\in\U$: an endpoint in $\D$ would leave some points of $\D$ in the
interval, because $\D$ has no greatest element. A bottom-end endpoint
would have the same problem. Now
\[
 y=\cut{\N}{a_i:i\in I}
\]
is positive infinite but below every $a_i$, so belongs to none of the
intervals, a contradiction. Hence $B^{-1}((\tfrac12,\tfrac32))=\U$ is
not open in this set-union sense, despite pointwise local continuity.

However, $F_1$ is an order automorphism by \cref{thm:auto}, so it \emph{is}
a homeomorphism for this topology: preimages of basic intervals are
intervals with the corresponding inverse-image endpoints, and preimages
preserve set-sized unions. The same is true of its inverse. Thus demanding
that the primitive itself be continuous in this global sense does not
remove the example. The original paper already describes the piecewise
$F_1$ as strongly continuous; this note uses the directly proved order
property rather than relying on that terminology.

There is no valid shortcut from continuity of $F_1$ and $\id$ to global
continuity of their difference in this particular framework. Indeed,
that difference is $-B$, for which the preceding inverse-image obstruction
applies. The familiar closure of continuous real-valued functions under
arithmetic must not be assumed without checking compatibility with the
chosen proper-class topology.

\subsection{Why Fornasiero's preservation theorem is not contradicted}
Fornasiero proves preservation of the sup property under specified uniform
recursive constructions with an additional \emph{tameness} hypothesis
\source{Theorem~3}{Fornasiero}. In his Definition~3.1, a nonconstant tame
function must, on each side of a gap and sufficiently near it, stay
strictly above or strictly below any specified comparison value.

Our $B$ fails that condition. At the finite--infinite gap, choose comparison
value $0$. On the entire left side $B=0$; no left neighborhood has
$B>0$ everywhere or $B<0$ everywhere. The global constant-function exception
does not apply because $B$ is nonconstant. Thus a published theorem about
tame functions cannot be used to infer the sup property for $B$.
Uniformity alone and tameness are different restrictions.

\subsection{A correct narrow uniqueness statement}
\begin{proposition}[Rational primitives remain unique]\label{prop:rational}
Let $K$ be an ordered field of characteristic zero. If two rational
functions without poles on a nonempty open interval have equal pointwise
derivatives there, then their difference is constant on that interval.
The same argument applies to rational surreal functions.
\end{proposition}
\begin{proof}
Write their difference as $P/Q$ in lowest terms over $K$, with $Q$ nonzero
on the interval. Differentiating algebraically gives
$(P'Q-PQ')/Q^2$. Since it vanishes at every point of the interval and the
interval contains infinitely many elements, the numerator polynomial is
identically zero.

If $P=0$ we are done. Otherwise, $P'Q=PQ'$ and $\gcd(P,Q)=1$ imply
$P\mid P'$. In characteristic zero a nonconstant $P$ has
$\deg P'<\deg P$, so this is impossible unless $P$ is constant.
Then the identity forces $Q'=0$, making $Q$ constant as well.
All operations involve only finitely many surreal coefficients, so the
same polynomial argument is legitimate over the proper-class field.
\end{proof}

Thus the countable \emph{cut} is doing real work. It converts signs of
rational functions into an infinitary Boolean condition; the resulting
function is not itself a rational function. A repair based on a smaller
function class must control this operation rather than only the
regularity of individual rational option functions.

\section{An ordered-field analogue and the hyperreal distinction}\label{sec:hyperreal}
\begin{proposition}\label{prop:field}
Let $K$ be any non-Archimedean ordered field. Define
\[
 b_K(x)=\begin{cases}
 1,&x>n\text{ for all ordinary }n\in\N,\\
 0,&\text{otherwise}.
 \end{cases}
\]
As an unrestricted function $K\to K$, it is nonconstant, invariant under
finite translations, and has order-field derivative zero everywhere.
\end{proposition}
\begin{proof}
Non-Archimedeanness supplies a positive element greater than every natural
number, whereas $b_K(0)=0$. The integer-bound proof of
\cref{lem:translation} uses only ordered-field arithmetic and applies in
$K$. In particular, $b_K(x+h)=b_K(x)$ for $|h|<1$, and every relevant
punctured difference quotient is exactly zero.
\end{proof}

The exceptional feature in the surreal problem is therefore not the mere
existence of such a function. It is the ability to certify it by the
published \emph{genetic syntax}. The Conway cut provides that certificate.

In a hyperreal setting with the usual internal/external distinction,
$b_K$ must be treated as an external function, not as the transferred
extension of an ordinary real function. Here is a direct diagnostic,
assuming the standard transfer framework. Were $b_K$ internal, the class
of hyperintegers $H$ with $b_K(H)=1$ would be an internal nonempty subset
of the hypernaturals. The transferred least-element principle would give
it a least member $H$. But $H$ is greater than every ordinary natural, and
so is $H-1$, contradicting minimality. This argument invokes the
least-element principle only for an internal set, which is precisely
where the contradiction enters.

The ordered-field derivative calculation is valid independently of such
extra structure. It neither refutes transferred theorems for internal
objects nor identifies genetic functions with internal hyperreal
functions. These are different restrictions on which functions are
admitted.

\section{Finite truncations and exact computational checks}\label{sec:checks}
\subsection{What finite option sets do}
Let
\[
 T_N(x)=\cut{q(n-x):0\le n\le N}{},\qquad B_N(x)=1-T_N(x).
\]
The maximum of the option-indexing set is $N$, so the binary cut lemma gives
\[
 T_N(x)=\begin{cases}1,&x\le N,\\0,&x>N.\end{cases}
\]
Each $B_N$ is an ordinary step at the actual surreal point $N$, where it
is not differentiable. In contrast, the countably indexed $B$ has no
exceptional point.

For every fixed surreal $x$, the sequence $T_N(x)$ is eventually exactly
$T(x)$. If $x\in\D$, choose a natural upper bound $n$ for $x$; all
$N\ge n$ give $T_N(x)=1$. If $x\in\U$, every $N$ gives $T_N(x)=0$.
This eventual stabilization is stronger than a numerical convergence
claim at the chosen input, but it is \emph{not uniform} in $x$: at
$x=N+1$, $T_N(x)=0$ while $T(x)=1$.

In particular, one must not replace a countable option family by a large
finite list and then assert that the resulting function has zero
derivative everywhere. The exact construction is infinitary, even though
its sign argument is short.

\subsection{The supplied verifier}
The file \texttt{code/verify.py} uses Python's standard library and exact
rational arithmetic. It implements the ordered rational-function field
$\Q(t)$ with $t$ positive infinite. For a nonzero rational function
$p(t)/r(t)$ with positively normalized denominator, its sign is determined
by the leading coefficient of $p$, and it is positive infinite exactly
when that sign is positive and $\deg p>\deg r$. This is a symbolic
classification, not a test against finitely many large integers.

The program separately enumerates a finite early-day dyadic surreal tree
to find the simplest number above selected finite sets of rational left
options. This checks the finite-cut implementation without hard-coding
the output of each example. It does not claim to implement arbitrary
surreal birthdays or countable cuts.

The archived run completed \textbf{7,248 exact checks}, all passing,
with a fixed random seed $20260920$ and $84$ rational-function inputs.
These are assertion counts, not counts of independently proved theorems.
The groups are summarized below; the machine-readable details are in
\texttt{results/verification.json}.

\begin{center}
\small
\begin{tabular}{lr}
\toprule
Check group & Assertions\\
\midrule
Rational-field identities &335\\
Sign and rational bounds &336\\
Polynomial square certificates &168\\
Explicit natural upper-bound witnesses &134\\
Selected infinite-option signs &85\\
Inverse maps &168\\
Punctured-neighborhood bounds &6\\
Local constancy &504\\
Primitive difference quotients &1008\\
Indicator difference quotients &504\\
Preservation of order &3486\\
Finite-translation group law &180\\
Finite Conway-cut searches &270\\
Triangular scale-evaluation matrix &64\\
\midrule
Total &7248\\
\bottomrule
\end{tabular}
\end{center}

These computations can expose arithmetic, sign, boundary-case, and
implementation mistakes. They do not verify membership in a published
genetic class, settle a quantifier over all surreal numbers, show that a
set-sized cut exists, or establish novelty. Those parts of the argument
are explicitly mathematical or bibliographic, and remain separate from
the test report.

\subsection{Reproduction}
From the archive directory, run
\begin{verbatim}
python code/verify.py
latexmk -pdf -interaction=nonstopmode -halt-on-error article.tex
latexmk -pdf -interaction=nonstopmode -halt-on-error short_proof.tex
\end{verbatim}
No network access or third-party Python packages are required for the
verification. The TeX documents use standard TeX Live packages; a build
script for PowerShell and a POSIX shell are supplied. External papers are
referenced, not redistributed.

\section{Conclusions and remaining questions}\label{sec:conclusion}
The central algebraic observation is that the cut
\[
 \cut{q(a-x):a\in A}{}
\]
turns a set-indexed family of sign tests into an existential assertion:
it equals $1$ exactly when $x\le a$ for at least one $a\in A$. For an
indexing set with no largest element, this assertion changes truth value
across an unrealized gap, not at any surreal point. A function can
therefore be locally constant at every input while being globally
nonconstant, and the broad genetic constructors admit it.

Under the definitions checked in this note, the outcomes are exact.
Genetic primitives need not be unique up to a constant; even real
normalization and an increasing-bijection requirement do not restore
uniqueness. Entire genetic functions need not have the sup property.
The pointwise differentiation kernel contains the ordinal-indexed
independent family of \cref{thm:independent}. These are consequences of
the explicit cut, not extrapolations from computation.

Several narrower research questions remain worth separating from the
resolved targets. What restrictions on option families preserve the sup
property? Which recursive subclasses exclude gap indicators while still
containing the desired transcendental functions? Can one characterize
an analytically useful class in which pointwise differentiation has only
constant kernel? Restricting to finitely many rational option expressions
removes the present schema, but no assertion that this restriction alone
settles those broader questions is made here.

The appropriate next external check is focused: verify the admissibility
certificate against the intended genetic-function conventions, and search
for prior publication of this countable-cut observation. The core proof
is short enough for such a check to be independent of the implementation
and of every later generalization in this article.

\appendix
\section{Proof dependency and objection audit}\label{app:audit}
\begin{longtable}{@{}p{0.29\textwidth}p{0.66\textwidth}@{}}
\toprule
Potential objection & Explicit resolution\\
\midrule
\endhead
The left options may form a proper class.
 & They are indexed by $\N$, or by a specified set $A$. For each input,
 replacement gives a set of values.\\[0.5em]
The cut may fail to define a number.
 & Its right option set is empty. Moreover every left option lies strictly
 below $1$, so even an explicit separating candidate is available.\\[0.5em]
The formula may hide a conditional option rule.
 & Every option $q(n-x)$ is always included and always defined.
 The finite/infinite split is derived from its sign, not used in the
 genetic definition.\\[0.5em]
The denominator may vanish at a nonreal input.
 & $1+t^2>0$ in every ordered field, including $\No$.\\[0.5em]
Only a nongenetic piecewise function has been constructed.
 & The nonrecursive option terms fit the published affine-template
 constructor, and the fixed option family is countable.\\[0.5em]
The construction may fail uniformity.
 & The expressions ignore the option variables of the argument. All
 representation choices leave exactly the same numerical cut.\\[0.5em]
The later cofinality condition may exclude it.
 & For every $y<x<z$, each left option is still $q(n-x)<T(x)$;
 right inequalities are vacuous.\\[0.5em]
A jump may occur at $\omega$.
 & The finite--infinite gap is not $\omega$. Every number in
 $(\omega-1,\omega+1)$ is positive infinite; $B$ is constant there.\\[0.5em]
The derivative may use only real increments.
 & The quotient identity holds for every surreal $h$ with $0<|h|<1$,
 including arbitrary infinitesimals.\\[0.5em]
A bounded interval might restore uniqueness.
 & $(0,\omega)$ already contains the witnesses $1$ and $\omega/2$.\\[0.5em]
The primitive might lack monotonicity.
 & $F_1=\Phi_{-1}$ is an increasing bijection, with inverse $\Phi_1$.\\[0.5em]
The sup property might allow the end gaps.
 & The witness class is nonempty and bounded by the surreal $\omega$.
 Neither end gap can be its least upper bound.\\[0.5em]
A known theorem about tame functions may contradict the result.
 & $B$ is not tame in Fornasiero's strict gap-comparison sense.\\[0.5em]
The finite computations might be presented as a formal proof.
 & Their scope is limited to exact finite shadows; the proper-class and
 countable-cut claims are proved in the text, not verified by the program.\\[0.5em]
The primitive itself may already be known.
 & It is: the source prints its piecewise form. Only the genetic
 realization and the deductions are presented as the contribution here.\\
\bottomrule
\end{longtable}

\section{Source locations and status of the search}\label{app:sources}
The primary-text locations checked for this note are as follows. Page
numbers are printed page numbers, not PDF-viewer indices.
\begin{center}
\small
\begin{tabularx}{\textwidth}{@{}lX@{}}
\toprule
Source & Locations used\\
\midrule
\cite{RSS} & Section~2.1 (cut and simplicity conventions);
 pp.~9--11 (genetic constructor, closure, rational unit step);
 pp.~27--28 (function limits and derivatives);
 p.~37, equation~(7) and Conjecture~53.\\[0.4em]
\cite{Berenbeim} & Definition~55, pp.~126--128 (adjunction and global
 cofinality); pp.~133--134 (entire functions and rational step);
 Chapter~8, pp.~225--226, Conjecture~3.\\[0.4em]
\cite{Fornasiero} & Example~2.5 (unused option variables);
 Definitions~3.1--3.2 and Theorem~3 (tameness and sup property).\\
\bottomrule
\end{tabularx}
\end{center}

Searches on 20 September 2026 included exact-conjecture queries and
combinations of ``genetic,'' ``surreal,'' ``primitive,'' ``counterexample,''
and ``sup property,'' together with the authors' names. Relevant results
located the original article, its arXiv version, the dissertation, and
Fornasiero's preprint. Some broader queries returned mostly irrelevant
material, so absence from these results is weak evidence about priority.
No located result supplied a later resolution of the two targets.
The source manifest records this limitation rather than labelling an
unsuccessful search as an exhaustive literature review.

The theorem statements and proofs in this note were developed from the
explicit formulas displayed here. Their evaluation does not depend on a
claim that the questions remained universally recognized as open on the
access date. If earlier work contains the same observation, the
mathematical counterexamples stand but the novelty claim must be revised.

\begin{thebibliography}{9}
\small
\bibitem{RSS}
Simon Rubinstein-Salzedo and Ashvin Swaminathan,
\emph{Analysis on Surreal Numbers},
Journal of Logic \& Analysis \textbf{6}:5 (2014), 1--39.
\href{https://doi.org/10.4115/jla.2014.6.5}{doi:10.4115/jla.2014.6.5}.
\href{https://math.hawaii.edu/home/jla/210-758-1-PB.pdf}{Published PDF};
\href{https://arxiv.org/abs/1307.7392}{arXiv:1307.7392}.

\bibitem{Berenbeim}
Alexander Michael Berenbeim,
\emph{Foundations for the Analysis of Surreal-Valued Genetic Functions},
Ph.D. dissertation, University of Illinois at Chicago, 2022.
\href{https://indigo.uic.edu/articles/thesis/Foundations_for_the_Analysis_of_Surreal-Valued_Genetic_Functions/20253705/1/files/36188862.pdf}{UIC repository PDF}, version~1, file~36188862.

\bibitem{Fornasiero}
Antongiulio Fornasiero,
\emph{Recursive definitions on surreal numbers},
arXiv:math/0612234, version~1, 9 December 2006.
\url{https://arxiv.org/abs/math/0612234}.
\end{thebibliography}
\end{document}
